\documentclass[12pt]{article}
\usepackage[a4paper,margin=1.25cm]{geometry}
\usepackage{amsmath,amssymb,mathtools}
\usepackage{graphicx}
\usepackage{booktabs}
\usepackage{tabularx}
\usepackage{tikz}
\usepackage[T1]{fontenc}
\usepackage[utf8]{inputenc}
\usepackage[english]{babel}

\setlength{\parindent}{0pt}
\setlength{\parskip}{0.45em}
\setlength{\abovedisplayskip}{5pt plus 1pt minus 1pt}
\setlength{\belowdisplayskip}{5pt plus 1pt minus 1pt}
\setlength{\abovedisplayshortskip}{3pt plus 1pt minus 1pt}
\setlength{\belowdisplayshortskip}{3pt plus 1pt minus 1pt}

\newcommand{\dd}{\,\mathrm{d}}
\newcommand{\e}{\mathrm{e}}

\begin{document}

\begin{center}
{\Large\bfseries A sequence pinned by its own partial sums}\\[0.35em]
{\large \(S_n=n^2f(n)\) forces \(f(n)=\dfrac{2f(1)}{n(n+1)}\)}
\end{center}

\textbf{Problem.} A function \(f\) on the positive integers satisfies
\[
\sum_{k=1}^{n}f(k)=n^2f(n)\qquad\text{for every }n\ge1,
\qquad f(1)=2027 .
\]
Find \(f(2026)\).

\[
\boxed{\displaystyle
f(n)=\frac{2f(1)}{n(n+1)},
\qquad
f(2026)=\frac{2\cdot2027}{2026\cdot2027}=\frac1{1013}.}
\]

\textbf{Overview.}
A condition relating a sequence to its own partial sums is really a
first-order recurrence in disguise. Subtracting the relation at \(n-1\)
from the relation at \(n\) eliminates the sum entirely and leaves
\[
(n+1)f(n)=(n-1)f(n-1),
\]
whose solution is a telescoping product collapsing to
\(\frac{2}{n(n+1)}\) times \(f(1)\). The value \(2027\) then cancels
against the \(2027\) hidden inside \(2026\cdot2027\), which is why the
answer is a clean \(\frac1{1013}\).

\section*{1. From the sum condition to a recurrence}

Write
\[
S_n=\sum_{k=1}^{n}f(k),
\qquad\text{so the hypothesis is}\qquad
S_n=n^2f(n)\quad (n\ge1).
\tag{1}
\]
For \(n=1\) this reads \(f(1)=1\cdot f(1)\), which is automatic: the
hypothesis places no constraint on \(f(1)\), and that is why the problem
must supply it.

For \(n\ge2\), subtract (1) at \(n-1\) from (1) at \(n\), using
\(S_n-S_{n-1}=f(n)\):
\[
n^2f(n)-(n-1)^2f(n-1)=f(n).
\]
Collecting the \(f(n)\) terms,
\[
(n^2-1)f(n)=(n-1)^2f(n-1),
\]
and factoring \(n^2-1=(n-1)(n+1)\), then cancelling the common factor
\(n-1\ne0\),
\[
\boxed{\displaystyle (n+1)f(n)=(n-1)f(n-1),\qquad n\ge2 .}
\tag{2}
\]
Equivalently
\[
\frac{f(n)}{f(n-1)}=\frac{n-1}{n+1},\qquad n\ge2 .
\tag{3}
\]
At \(n=2\) relation (2) reads \(3f(2)=1\cdot f(1)\), so
\(f(2)=\frac{f(1)}{3}\); the recurrence is well posed from the start,
and the sum has disappeared completely.

\section*{2. Solving the recurrence}

Iterating (3) from \(n\) down to \(2\),
\[
f(n)=f(1)\prod_{k=2}^{n}\frac{k-1}{k+1}
=f(1)\cdot\frac{1\cdot2\cdot3\cdots(n-1)}{3\cdot4\cdot5\cdots(n+1)} .
\tag{4}
\]
The numerator is \((n-1)!\) and the denominator is
\(\frac{(n+1)!}{2}\), so the product collapses:
\[
\prod_{k=2}^{n}\frac{k-1}{k+1}
=\frac{(n-1)!}{(n+1)!/2}
=\frac{2}{n(n+1)} .
\]
Hence
\[
\boxed{\displaystyle f(n)=\frac{2f(1)}{n(n+1)},\qquad n\ge1,}
\tag{5}
\]
and the formula is also correct at \(n=1\), where it returns
\(\frac{2f(1)}{2}=f(1)\).

The mechanism is a shifted telescope: the numerator of each factor is the
denominator of the factor two steps earlier, so all but two terms cancel
at each end, one pair surviving at the bottom (\(1\cdot2\)) and one
pair at the top (\(n(n+1)\)).

\section*{3. Verification}

Formula (5) must still be checked against (1), not merely against (2),
since the passage to the recurrence discards information. With
\(f(k)=\frac{2f(1)}{k(k+1)}\),
\[
S_n=2f(1)\sum_{k=1}^{n}\frac1{k(k+1)}
=2f(1)\sum_{k=1}^{n}\left(\frac1k-\frac1{k+1}\right)
=2f(1)\left(1-\frac1{n+1}\right)
=\frac{2f(1)\,n}{n+1},
\]
while
\[
n^2f(n)=n^2\cdot\frac{2f(1)}{n(n+1)}=\frac{2f(1)\,n}{n+1}.
\]
The two agree for all \(n\ge1\), so (5) is the unique solution. \(\square\)

The partial-fraction telescope \(\frac1{k(k+1)}=\frac1k-\frac1{k+1}\) is
the same collapse as in Section 2, now on the additive side, a sign
that \(n^2\) was not an arbitrary weight.

\section*{4. The answer}

With \(f(1)=2027\) and \(n=2026\),
\[
f(2026)=\frac{2\cdot2027}{2026\cdot2027}
=\frac{2}{2026}
=\frac1{1013}.
\]
The constant \(2027\) cancels identically: it appears in the numerator
through \(f(1)\) and in the denominator as the factor \(n+1=2027\). The
problem is arranged so that the requested index is exactly one below the
given initial value, which is the only reason the answer is rational with
such a small denominator. Since \(2026=2\cdot1013\), the final form is
\(\frac1{1013}\).

\begin{center}
\small
\renewcommand{\arraystretch}{1.3}
\begin{tabularx}{0.8\textwidth}{@{}lXXX@{}}
\toprule
\(n\) & \(f(n)=\frac{2\cdot2027}{n(n+1)}\) & \(S_n\) & \(n^2f(n)\) \\
\midrule
\(1\) & \(2027\)                & \(2027\)                  & \(2027\) \\
\(2\) & \(2027/3\)              & \(4054/3\)                & \(4054/3\) \\
\(3\) & \(2027/6\)              & \(3\cdot2027/2\)          & \(3\cdot2027/2\) \\
\(4\) & \(2027/10\)             & \(8\cdot2027/5\)          & \(8\cdot2027/5\) \\
\bottomrule
\end{tabularx}
\end{center}

\section*{5. The general weight}

Nothing is special about \(n^2\) except the arithmetic that follows.
Suppose instead
\[
\sum_{k=1}^{n}f(k)=w(n)f(n),\qquad n\ge1 ,
\]
for a given sequence \(w\). Consistency at \(n=1\) forces \(w(1)=1\).
The same subtraction gives
\[
\bigl(w(n)-1\bigr)f(n)=w(n-1)f(n-1),
\qquad\text{so}\qquad
f(n)=f(1)\prod_{j=2}^{n}\frac{w(j-1)}{w(j)-1},
\tag{6}
\]
valid whenever \(w(n)\ne1\) for \(n\ge2\).

\begin{center}
\small
\renewcommand{\arraystretch}{1.5}
\begin{tabularx}{\textwidth}{@{}>{\raggedright\arraybackslash}p{0.14\textwidth}>{\raggedright\arraybackslash}p{0.26\textwidth}X@{}}
\toprule
\(w(n)\) & ratio \(\dfrac{w(n-1)}{w(n)-1}\) & solution \\
\midrule
\(n\)
& \(\dfrac{n-1}{n-1}=1\)
& \(f(n)=f(1)\): the condition \(S_n=n f(n)\) says every term equals the
running average, i.e. \(f\) is constant. \\
\addlinespace
\(n^2\)
& \(\dfrac{(n-1)^2}{n^2-1}=\dfrac{n-1}{n+1}\)
& \(f(n)=\dfrac{2f(1)}{n(n+1)}\), the present problem. \\
\addlinespace
\(2n-1\)
& \(\dfrac{2n-3}{2n-2}\)
& \(f(n)=f(1)\displaystyle\prod_{j=2}^{n}\frac{2j-3}{2j-2}
=f(1)\frac{(2n-3)!!}{(2n-2)!!}\sim\dfrac{Cf(1)}{\sqrt n}\). \\
\addlinespace
\(n^3\)
& \(\dfrac{(n-1)^2}{n^2+n+1}\)
& no elementary closed form; the denominator \(n^2+n+1\) does not factor
over \(\mathbb Q\), so the product does not telescope. \\
\bottomrule
\end{tabularx}
\end{center}

The table isolates what makes the stated problem work: for \(w(n)=n^2\)
the difference \(w(n)-1=n^2-1\) factors, and one of its factors cancels
the \((n-1)^2\) coming from \(w(n-1)\). A weight is ``nice'' for this
problem exactly when \(w(n)-1\) shares a factor with \(w(n-1)\), and
\(n^2-1=(n-1)(n+1)\) is the smallest interesting instance.

\section*{6. A continuous shadow}

The same computation has an integral analogue that explains the exponent.
If \(F(x)=\int_1^x f\) satisfies \(F(x)=x^2f(x)=x^2F'(x)\), then
\[
\frac{F'}{F}=\frac1{x^2},
\qquad
\log F(x)=-\frac1x+C,
\qquad
F(x)=A\,\e^{-1/x},
\]
so \(f(x)=F'(x)=\frac{A}{x^2}\e^{-1/x}\), which decays like \(x^{-2}\).
The discrete solution \(f(n)=\frac{2f(1)}{n(n+1)}\) has the same
\(n^{-2}\) decay; the discrete problem simply replaces the exponential
factor by an exact rational one, because the discrete telescope is exact
where the continuous one requires an integrating factor.

\end{document}
